\documentclass[a4paper,11pt]{article}
\usepackage[T1]{fontenc}
\usepackage[utf8]{inputenc}
\usepackage[italian]{babel}
\usepackage{amsmath,amssymb}
\usepackage{geometry}
\usepackage{xcolor}
\usepackage{listings}
\geometry{margin=2.5cm}

\lstdefinelanguage{MiniZinc}{
  morekeywords={int,set,of,array,var,constraint,include,solve,satisfy,output,show},
  sensitive=true,
  morecomment=[l]{\%},
  morestring=[b]"
}

\lstset{
  language=MiniZinc,
  basicstyle=\ttfamily\small,
  keywordstyle=\color{blue}\bfseries,
  commentstyle=\color{teal!70!black}\itshape,
  stringstyle=\color{orange!70!black},
  numbers=left,
  numberstyle=\tiny\color{gray},
  stepnumber=1,
  numbersep=8pt,
  showstringspaces=false,
  frame=single,
  breaklines=true
}

\title{CSP: n-Queens}
\date{}

\begin{document}
\maketitle

Si vuole trovare, per un dato intero $n>3$, una disposizione di $n$ regine su una
scacchiera $n\times n$ tale che nessuna coppia di regine sia in conflitto (stessa riga,
stessa colonna, stessa diagonale).

\subsection*{1.1 Modellazione}

Una formulazione CSP ad alto livello è la seguente.

\paragraph{Variabili}
Definiamo una variabile per ogni colonna:
\[
X = \{x_1, x_2, \dots, x_n\},
\]
dove $x_i$ rappresenta la riga in cui è posizionata la regina nella colonna $i$.

\paragraph{Domini}
\[
\forall i \in \{1,\dots,n\}:\quad D(x_i)=\{1,\dots,n\}.
\]

\paragraph{Vincoli}
\[
\begin{aligned}
C=\{\,&\text{alldifferent}(x_1,\dots,x_n),\\
      &\text{alldifferent}(x_1+1,\dots,x_n+n),\\
      &\text{alldifferent}(x_1-1,\dots,x_n-n)\,\}.
\end{aligned}
\]

Il primo vincolo impone che due regine non stiano sulla stessa riga.

Il vincolo di colonna è già soddisfatto implicitamente (per costruzione).

I vincoli sulle diagonali sono modellati con due \texttt{alldifferent}:
Due regine nelle posizioni $(i,x_i)$ e $(j,x_j)$ sono sulla stessa diagonale
principale se e solo se $x_i-i=x_j-j$, e sulla stessa diagonale secondaria se e solo se
$x_i+i=x_j+j$. 
Quindi imporre \texttt{alldifferent} su tutti i valori $x_i-i$ e $x_i+i$
elimina entrambi i tipi di conflitto diagonale.

\subsection*{1.2 Istanziazione per $n=5$}

Si ottiene il CSP $(X,D,C)$:

\[
X=\{x_1,x_2,x_3,x_4,x_5\}.
\]

\[
D(x_1)=D(x_2)=D(x_3)=D(x_4)=D(x_5)=\{1,2,3,4,5\}.
\]

\[
\begin{aligned}
C=\{\,&\text{alldifferent}(x_1,x_2,x_3,x_4,x_5),\\
      &\text{alldifferent}(x_1+1,x_2+2,x_3+3,x_4+4,x_5+5),\\
      &\text{alldifferent}(x_1-1,x_2-2,x_3-3,x_4-4,x_5-5)\,\}.
\end{aligned}
\]

\newpage

\subsection*{1.3 Codifica in MiniZinc}

\begin{lstlisting}
include "alldifferent.mzn";
int: n;
array[1..n] of var 1..n: q;

constraint alldifferent(q);
constraint alldifferent([q[i] + i | i in 1..n]);
constraint alldifferent([q[i] - i | i in 1..n]);

solve satisfy;

output [show(q)];
\end{lstlisting}


\end{document}